\section{Implementasi}

\subsection{Lingkungan dan Dependensi}
% TODO: Elixir/OTP, library (OTP \texttt{crypto} dan \texttt{public\_key}), versi, tautan dokumentasi.

\subsection{Implementasi Tiap Algoritma}
% TODO: parameter dan pemanggilan API untuk RSA-PSS, ECDSA (konversi DER ke $r \mid s$), Ed25519.

\subsection{Antarmuka CLI}
% TODO: subcommand dan contoh pemakaian.

\subsection{Distribusi Container}
% TODO: Dockerfile, base image, multi-arch, digest, pemetaan versi program, tag repositori, dan digest image.

\subsection{Integrasi CI/CD}
% TODO: workflow build image, workflow pemakai (Anmategra), pengelolaan secret dan trust store.

\subsection{Subcommand SBOM}
Pembungkusan dilakukan oleh \texttt{sbom bind}, yang membaca artefak dan berkas SBOM, memastikan
SBOM berupa objek JSON dan tidak lebih dari 16 MiB, lalu menuliskan pembungkus. Penulisan tidak
menimpa berkas yang sudah ada. \texttt{sbom extract} mengembalikan byte SBOM asli dari
pembungkus tanpa perubahan, sehingga kakas hilir dapat membaca berkas yang sama dengan keluaran
kakas asalnya.

\begin{table}[H]
\centering
\begin{tabular}{@{}p{2.8cm}p{5.6cm}p{5cm}@{}}
\toprule
Subcommand & Parameter & Keluaran \\
\midrule
\texttt{sbom bind} & \texttt{-file}, \texttt{-sbom}, \texttt{-out} &
Pembungkus JSON pada \texttt{-out}. \\
\texttt{sbom extract} & \texttt{-sbom}, \texttt{-out} &
Byte SBOM asli pada \texttt{-out}. \\
\texttt{sbom sign} & \texttt{-sbom}, \texttt{-package}, \texttt{-version}, \texttt{-key}, \texttt{-out} &
Bundle pembungkus pada \texttt{-out}. \\
\texttt{sbom verify} & \texttt{-file}, \texttt{-bundle}, \texttt{-sbom}, \texttt{-sbom-bundle}, \texttt{-trust-store}, \texttt{-package}, \texttt{-version}, \texttt{-json} (opsional) &
Hasil verifikasi sesuai Tabel 2.9. Exit code sesuai Tabel 2.10. \\
\bottomrule
\end{tabular}
\caption{Subcommand tambahan untuk SBOM.}
\end{table}

\subsubsection*{Pengkodean Base64 pada pembungkus}
Byte SBOM disimpan di dalam pembungkus dalam bentuk Base64 standar, bukan sebagai objek JSON
bersarang. Ada dua alasan utama.

\begin{enumerate}
  \item \textbf{Byte SBOM tetap utuh.} Tanda tangan dibuat atas byte pembungkus, sehingga
    byte SBOM di dalamnya juga harus dijaga. Jika SBOM ditulis sebagai objek JSON bersarang,
    parser akan membacanya lalu menulisnya ulang. Spasi, urutan key, dan escape dapat berubah
    dalam proses itu. Akibatnya \texttt{extract} tidak lagi identik dengan keluaran kakasnya.
    Dengan Base64, \texttt{extract} mengembalikan byte asli yang diuji dengan \texttt{cmp}.
  \item \textbf{JSON tidak membawa byte sembarang dengan aman.} String JSON harus berupa UTF-8,
    dan beberapa karakter wajib di-escape. Base64 hanya memuat huruf, angka, \texttt{+}, \texttt{/},
    dan \texttt{=}, sehingga aman disisipkan sebagai string tanpa escape.
\end{enumerate}

Pilihan ini juga menghindari masalah parser yang ketat. Kakas seperti Trivy dapat menulis
\texttt{null}, yang ditolak oleh \texttt{StrictJSON}. Karena isi SBOM tidak diparse oleh
\texttt{StrictJSON}, \texttt{null} di dalamnya tidak menjadi masalah. Program hanya memeriksa
bahwa isinya objek JSON, memakai \texttt{JSON.decode} bawaan.

Pilihan ini memiliki \textit{trade-off}. Pertama, ukuran naik sekitar sepertiga akibat Base64,
sehingga batas SBOM ditetapkan 16 MiB dan ukuran pembungkus dihitung dari batas tersebut.
Kedua, pembungkus tidak dapat dibaca langsung oleh manusia dan perlu didekode dengan
\texttt{sbom extract} untuk melihat isi SBOM. Ketiga, program tidak memeriksa isi SBOM lebih
jauh daripada bentuk objek JSON, sehingga kelengkapan komponen menjadi tanggung jawab kakas
pembuat SBOM.

Contoh alur dengan Trivy:

\begin{lstlisting}[basicstyle=\ttfamily\footnotesize,caption={Alur SBOM dengan kakas eksternal},label={lst:sbom-flow},numbers=none]
trivy fs --format cyclonedx --scanners secret --output sbom.json dist/
nadina sbom bind --file dist/contoh-app.zip --sbom sbom.json --out pembungkus.json
nadina sbom sign --sbom pembungkus.json --package contoh-app --version 1.0.0 \
  --key kunci.key.pem --out sbom-bundle.json
nadina sbom verify --file dist/contoh-app.zip --bundle rilis.json --sbom pembungkus.json \
  --sbom-bundle sbom-bundle.json --trust-store trust.json \
  --package contoh-app --version 1.0.0 --json
\end{lstlisting}
